% =====================================================================
%  Conferenze, Zagarolo (On The Road), giorno 1 — Intervista 3
%  Perché l'industria italiana fatica e che cosa possiamo fare insieme
%  Fonte: trascrizione automatica "Track 11". Intervista pomeridiana a
%  Davide, rappresentante dell'associazione di categoria dell'industria
%  italiana dei videogiochi (IIDEA), condotta dal moderatore
%  dell'evento, con molte domande dal pubblico. Audio spesso disturbato.
%  Etichette con prefisso "c2-3:". Nessun pacchetto aggiuntivo.
% =====================================================================

\chapter[Perché l'industria italiana fatica]{Perché l'industria italiana fatica,
e che cosa possiamo fare insieme}
\label{cap:c2-3-industria-italiana}

L'ospite del pomeriggio del primo giorno è, dice il moderatore, una persona che
``per gli studi in Italia non è non nota: lo conoscono praticamente tutti''. È
\textbf{Davide}\footnote{Il cognome non è pronunciato nella parte registrata. Con
ogni probabilità è la stessa persona che compare nel panel del
capitolo~\ref{cap:c1-1-tutela-minori}, dove parla a nome dell'associazione di
categoria e presenta lo stesso sito sui videogiochi italiani.}, che lavora per
\textbf{IIDEA}\footnote{Trascritto ``all'Idea'' e ``Idea''. IIDEA (Italian
Interactive Digital Entertainment Association) è l'associazione di categoria
dell'industria dei videogiochi in Italia.}, l'associazione di categoria
dell'industria italiana dei videogiochi. Il moderatore propone di fargli ``tutte le
domande che non si potrebbero fare in un'altra sede''; Davide ribatte che non si
tira mai indietro, nemmeno in contesti istituzionali, e paragona l'incontro a
``una bellissima seduta collettiva'' di psicoterapia.

% ---------------------------------------------------------------------
\section{Chi è e che cosa fa}
\label{sec:c2-3-ospite}

Davide apre scherzando su \emph{GTA VI}, giocato fino alle 2:30 di notte in albergo
in \emph{remote play} su iPad, e osserva una cosa ``bella'' dell'industria: tante
persone preferiscono stare \textbf{in presenza} a fare le cose insieme invece di
restare a casa a giocare.

In associazione fa \textbf{da tramite}: tra le aziende che sviluppano videogiochi e
l'associazione, ma anche nel \textbf{dialogo con le istituzioni}, con il resto
dell'industria e con gli enti internazionali. ``Sono sempre quello che sbaglia,
perché sto in mezzo'': l'associazione sta in mezzo a tante cose, e lui ``a maggior
ragione''. Prima ha lavorato \textbf{dieci anni nell'industria} con molti
cappelli: comunicazione, marketing, PR\footnote{Trascritto ``VR''; dal contesto è
più probabile ``PR''.}, giornalismo e, ``cosa che molti non sanno'',
\textbf{localizzazione}. ``Mi manca solo sviluppare giochi.''

Una premessa ``psicologica'': chi fa i videogiochi per lavoro è spesso un
\textbf{fanatico}. A chi ama il cinema non viene per forza voglia di lavorare nel
cinema, a chi ama la musica di suonare; a chi ama i videogiochi, invece, prima o
poi viene da pensare ``questo è il gioco che vorrei tanto fare''.

% ---------------------------------------------------------------------
\section{Perché gli studi italiani faticano}
\label{sec:c2-3-cause}

Prima domanda: dalla sua posizione privilegiata, perché gli studi italiani
faticano ad affermarsi?

\subsection{Il contesto globale}

Prima di tutto, un dato globale da ``mettersi in tasca'': \textbf{esistono troppi
videogiochi}, e troppa gente vuole farli, molti più di quanti se ne possano
giocare e scoprire. Spesso si guardano solo i successi degli altri: ma parlando con
i colleghi in Francia, in Spagna o altrove, nessuno dice ``siamo a posto''. Accanto
ai casi di successo c'è ovunque ``uno stuolo di aziende e di team che arrancano
allo stesso modo''.

\subsection{Le differenze italiane}

\begin{enumerate}
  \item \textbf{Densità dell'industria.} L'Italia è partita più tardi. All'inizio
        c'erano solo avventure ``più o meno colossali di pazzi visionari''; della
        prima ondata sono sopravvissute poche realtà che fanno ancora oggi ciò che
        facevano vent'anni fa (cita Milestone e pochi altri\footnote{La
        trascrizione nomina anche ``Ivano Venturi'': con ogni probabilità un
        riferimento a uno sviluppatore storico italiano, non verificabile.}).
        L'Italia è arrivata tardi anche al \textbf{riconoscimento} dell'industria e
        a qualsiasi forma di \textbf{sostegno}. Altrove, anche con popolazioni
        minori (il Nord Europa), lavorano molte più aziende e persone, e per la
        legge dei grandi numeri ci sono più successi.
  \item \textbf{Fare un gioco non è fare uno studio.} Una cosa è dire ``voglio fare
        questo videogioco'', un'altra ``voglio avere uno studio di sviluppo''.
        Spesso si viene dalla passione e si mescolano le due cose. Cita il caso di
        un gioco di qualche anno fa, ``otto anni \emph{in the making}'', fatto la
        sera da un artista visivo con famiglia e un altro lavoro\footnote{Il titolo
        è illeggibile nella trascrizione.}: un successo, ma di chi non voleva fare
        videogiochi per mestiere. Chi vuole fare impresa, invece, ha i problemi di
        \textbf{tutti gli imprenditori italiani}: \textbf{accesso al capitale} e
        \textbf{sostenibilità}.
  \item \textbf{Capitali difficili da attrarre.} Nei videogiochi l'accesso al
        capitale è ancora più complicato: in Italia ci sono ``troppe cose che vanno
        tendenzialmente bene'' e in cui si fanno soldi, perciò il capitale non ha
        bisogno di rischiare su un settore nuovo (a differenza, sembra intendere,
        di paesi come la Polonia, che nel videogioco hanno trovato un motore di
        sviluppo\footnote{Il passaggio è molto disturbato (``la Bologna, the
        Witcher''); il riferimento a \emph{The Witcher} fa pensare alla Polonia.}).
  \item \textbf{Supporto istituzionale arrivato tardi e ancora scarso.}
        L'associazione si è spesa molto per ottenere il \textbf{tax credit}:
        ``funziona, ma non è abbastanza''. La \textbf{legge cinema} prevederebbe
        anche \textbf{contributi selettivi} e \textbf{contributi automatici} per i
        videogiochi, che però non sono ancora operativi. ``Uno dei valori
        dell'associazione è rompere le scatole perché queste cose succedano.''
  \item \textbf{Mancanza di alcune competenze.} Su Steam ci sono tra i
        \textbf{500 e i 600} giochi italiani pubblicati. Il \textbf{42\%} ha
        \textbf{meno di 50 recensioni}. Eppure, negli ultimi tre anni, con
        un'impennata di uscite, il gradimento medio del pubblico è
        \textbf{superiore all'80\%} (recensioni positive, ``non quella disgrazia di
        Metacritic''). Quindi la \textbf{qualità sta migliorando}: c'è talento e
        competenza. Il problema è la \textbf{scoperta} (\emph{discoverability}):
        probabilmente si arriva all'uscita \textbf{senza più soldi}, e il gioco
        bello non trova il suo pubblico.
  \item \textbf{Dimensioni.} Il censimento dell'associazione conta circa
        \textbf{200 aziende} e \textbf{3.800 persone}; aggiungendo le realtà
        informali e quelle non censite il numero cresce, ma resta basso rispetto ad
        altri paesi, in una congiuntura di mercato ``non esattamente felice''.
\end{enumerate}

% ---------------------------------------------------------------------
\section{Come aiutare l'associazione a fare pressione}
\label{sec:c2-3-pressione}

Il moderatore toglie ``un velo'': molti pensano che la pressione
dell'associazione sullo Stato non sia abbastanza, rispetto alle agevolazioni che
esistono in Francia o in Germania. C'è un modo concreto in cui studi piccoli e
grandi e i tanti studenti presenti possono aiutare?

\subsection{Fare gruppo, non gruppi}

Davide comincia da ciò che \textbf{non} fare. In Italia c'è la tendenza a
\textbf{non fare gruppo, ma fare gruppi}: gruppetti disomogenei che restano uno
contro l'altro. ``Quando siamo quattro gatti'', farsi la guerra fa solo male al
settore. C'è anche confusione su \textbf{chi si occupa di che cosa}; e se non tutti
lo sanno, ammette, chi fa le cose ``dovrebbe comunicare meglio, anzi
sicuramente''.

Sul tax credit e sugli incentivi vede spesso l'idea che esista una ``lobby dei
grandi industriali'' che si spartisce ``una torta infinita di soldi''. Non è così:
tutto è \textbf{pubblico e tracciabile}. Ciò che si può fare:
\begin{itemize}
  \item \textbf{sapere come funzionano le cose}: che cosa fa un'associazione di
        categoria (un ente vicino all'industria ma di forma
        \textbf{semi-istituzionale}, che deve rispettare certi parametri), chi dà
        i soldi ai videogiochi in Italia, chi dovrebbe darli;
  \item \textbf{chiedere} quando ci sono dubbi: ``le risposte si hanno'';
  \item \textbf{sostenere le iniziative} comuni. Cita una campagna
        sul tax credit a giugno: ``magicamente, dopo tre giorni'' è arrivata una
        risposta. In un periodo di equilibri politici sensibili, fare pressione
        \textbf{come sistema} funziona.
\end{itemize}
Ma bisogna anche ``rendersi conto di chi siamo'': ci sono industrie molto più
numerose che, se fanno pressione, ``bloccano le strade''. Esercitare lo
\textbf{scetticismo}, sì, ma non il \textbf{sospetto}: ``c'è differenza tra
scetticismo, curiosità e sospetto''. Poiché la torta è piccola, l'istinto è
prendersi una briciola; invece bisogna capire insieme come gestire l'accesso.

\subsection{Non siamo ancora un'industria}

In Italia non c'è un'``industria alta'' e un'``industria bassa'': anche la più
grande azienda del settore, per l'ordinamento italiano, è poco più di una piccola
impresa. ``Ci raccontiamo come industria, ma \textbf{come numeri non siamo
un'industria}.'' L'unico modo per diventarlo è stare dalla stessa parte: se uno
studio fa un successo enorme, ``è per tutti''. ``Venite a rompermi le scatole'':
quella mattina, a chi gli chiedeva come fare una certa cosa, ha risposto ``non lo
so, ma scrivimi e ti faccio rispondere da chi se ne occupa''.

% ---------------------------------------------------------------------
\section{Interventi concreti}
\label{sec:c2-3-interventi}

Quali interventi concreti, come collettività, per far crescere l'industria?

\begin{itemize}
  \item \textbf{Capire che cosa si fa.} Se si vogliono fare videogiochi per lavoro,
        bisogna usare gli strumenti per interpretare il mercato; \textbf{accollarsi
        i fallimenti}, \textbf{fallire presto, fallire in fretta} e ripartire.
  \item \textbf{Pressione anche sul livello locale}: regioni e comuni. Racconta di
        aver visto il pitch di un gioco poi diventato famosissimo (``che canta un
        po' in francese''\footnote{Il titolo non è detto; il riferimento a un
        grande successo francese con un publisher alle spalle è chiaro, ma
        l'identificazione resta un'ipotesi.}): chi si siede al tavolo di un
        publisher dicendo ``ho queste agevolazioni fiscali, in questa regione
        prendo questo, in quest'altra quest'altro'' è \textbf{molto attrattivo},
        perché si porta ``una dote''.
  \item \textbf{Soldi all'inizio e alla fine.} In Italia mancano i soldi
        \textbf{all'inizio} (difficilissimo per tutti) e \textbf{alla fine} dello
        sviluppo, per vendere davvero il gioco in un mondo in cui ``l'attenzione
        dura un secondo''. Servono più iniziative che diano \textbf{piccole cifre}
        per fare un prototipo senza ``rischiare la vita'', e poi i contributi
        selettivi e automatici per chi non ha i numeri per accedere al tax
        credit (``in Italia funziona così: capitale minimo e via
        dicendo'').
  \item \textbf{Cicli più brevi.} Come detto al mattino
        (capitolo~\ref{cap:c2-1-marketing-publisher}), per il primo gioco
        \textbf{andare sul piccolo}: oggi bisogna arrivare presto sul mercato.
        Davide racconta di aver ``litigato su un palco'' con chi sostiene ``se hai
        quell'idea, vai dritto''. La sua posizione: l'obiettivo di chi fa
        videogiochi \textbf{non è fare il primo gioco, ma il secondo}. Se c'è
        un'idea grande ci si deve arrivare, ma non è detto al primo tentativo.
        ``Non è che il gioco non sia bello: è che lo giocano i vostri familiari'':
        brutto da dire, ma bisogna esserne consapevoli.
\end{itemize}

% ---------------------------------------------------------------------
\section{Le domande del pubblico}
\label{sec:c2-3-domande}

\subsection{Le istituzioni in vent'anni}

Come hanno reso difficile il lavoro le istituzioni, dal 2006 a oggi? Davide
individua un problema di \textbf{status medio} di chi occupa le cariche: persone
ancorate a un sistema produttivo tradizionale che funziona. Convincere qualcuno
che un piccolo videogioco possa essere ``una cosa concreta'' in Italia è molto
difficile.

Il problema più grande è culturale: un \textbf{pregiudizio altissimo} sul fatto che
i videogiochi non siano cultura. ``Mi fa impazzire'': i videogiochi \textbf{sono}
cultura proprio perché contengono una complessità di competenze, talenti,
conoscenze e \textbf{artigianalità} (``a volte devi reinventare la ruota''). Un
film è considerato cultura di per sé; per il videogioco ci si chiede ``che
videogioco è?''.

Ma le cose sono cambiate molto: le istituzioni \textbf{ascoltano} molto di più
(alcune fanno, altre ascoltano e dimenticano). L'evento stesso nasce da una
volontà istituzionale locale; a Roma ci sono state iniziative. Anche l'opinione
pubblica è cambiata: oggi i titoli sui giornali allarmistici su \emph{GTA VI} si
prendono ``quasi automaticamente sul ridere''. Il vero problema è la
\textbf{velocità}: le istituzioni non possono seguire i tempi dell'industria.

\subsection{Chi dovrebbe insegnare il mestiere}

Una sviluppatrice racconta un'esperienza con un \textbf{bando pubblico} a cui il suo
team (al primo gioco, con un budget comunale) ha partecipato: dopo aver perso il
programmatore, una persona coinvolta nell'organizzazione dell'iniziativa ha offerto
loro \textbf{tre mesi di un programmatore per 11.500 euro}, specificando che non era
il momento di fare richieste alte perché sapeva che avevano solo i soldi del comune.
Non è un altro problema di sistema, chiede: potersi fidare di chi dovrebbe insegnare
il mestiere?

``Assolutamente sì.'' È ciò che intende con ``parliamoci'': questi problemi vanno
\textbf{segnalati}, anche ufficialmente. L'associazione ha vincoli su quando può
intervenire, ma il fatto che i soldi siano pochi non significa che ``valga
tutto''. ``Fate i compiti a casa, non vi fate fregare; se ci sono problemi,
ditelo.''

\subsection{Quali competenze mancano}

Davide precisa: non mancano le \textbf{competenze tecniche}, che ci sono, e c'è
capacità di organizzazione (``siamo abituati a fare cattedrali con i
legnetti''\footnote{Trascritto ``le cattedrali con i regolotti''.}). Manca spesso la
\textbf{visione editoriale}: si parte dal progetto e dall'idea, non dal mercato e
dal pubblico; e mancano budget e competenze per \textbf{arrivare sul mercato}.

Aggiunge un punto a cui tiene: il fatto che l'Italia \textbf{costi poco} viene
spesso presentato come vantaggio competitivo, ed effettivamente lo è per attrarre
fondi, ma è \textbf{un'arma a doppio taglio}. Stipendi bloccati significano
incapacità di \textbf{trattenere il talento}: un senior a cui uno studio
internazionale offre il doppio, magari lavorando da casa, ``perché non dovrebbe
andare all'estero?''. Le aziende italiane hanno quindi un doppio problema: accedere
a capitali maggiori e alzare gli stipendi in un sistema che fa crescere costi e
tasse.

\subsection{Troppi sviluppatori, pochi amministratori}

Un'altra spettatrice osserva di vedere intorno a sé tantissimi sviluppatori e
\textbf{zero amministratori}: c'è un distacco tra fare un gioco e fare impresa.
Ringrazia la regione per i percorsi formativi e nota che, sfruttando bene le
opportunità esistenti (anche quelle legate alla disoccupazione), si potrebbe
arrivare a una prima produzione; ma manca, a livello nazionale, la diffusione di
queste conoscenze, e manca una forza politica chiaramente schierata a sostegno dei
videogiochi.

Davide è d'accordo: quando parla di ``convincere i distretti locali'' non intende
``piogge di soldi a caso'', ma anche \textbf{percorsi di formazione
imprenditoriale} specifici, che regioni e comuni possono offrire prima ancora dei
finanziamenti. La combinazione di \textbf{richiesta dal basso} e capacità
dell'associazione di formulare offerte che interessino le istituzioni funziona. ``Non
siamo lo sportello lamentele del mondo'', e in associazione sono ``cinque
persone'', ma ci provano.

\subsection{Cinque punti per il successo}

Un ultimo spettatore chiede ``cinque punti'' operativi. Davide li elenca, in sintesi:
\begin{enumerate}
  \item \textbf{Sostegno}: colmare il gap con le altre industrie europee; è
        fattibile ed è un obiettivo condiviso.
  \item \textbf{Abbassare la soglia del sospetto}.
  \item \textbf{Fare meglio i compiti a casa}, tutti, anche andando a ``rompere le
        scatole'' agli sviluppatori più importanti: in Italia, oggi, ``tutti
        rispondono a tutti''.
  \item \textbf{Chiedere come funzionano le cose}: alle istituzioni, all'associazione,
        a chiunque.
  \item \textbf{Ascoltare il mercato che cambia e agire}: fare i compiti non
        significa copiare ciò che funziona oggi, ma \textbf{capire come funzionano
        le cose} e farle, ``senza la sfera di cristallo''.
\end{enumerate}
In sintesi: le istituzioni e l'associazione devono garantire più accesso; l'industria
deve capire il mercato e adeguarsi, \textbf{senza perdere la propria identità}, che
è d'impresa ma anche di \textbf{espressione artistica}.

% ---------------------------------------------------------------------
\section{Il sito dei videogiochi italiani}
\label{sec:c2-3-sito}

Davide chiude presentando la nuova versione del sito che raccoglie
\textbf{tutti i videogiochi italiani}\footnote{L'indirizzo è trascritto in modo
illeggibile (``insignita di punto.it''); è lo stesso sito citato nel
capitolo~\ref{cap:c1-1-tutela-minori}.}, aggiornato quasi in tempo reale a partire
dai dati di Steam:
\begin{itemize}
  \item una \textbf{mappa degli studi} (``se non ci siete, iscrivetevi'');
  \item schede delle aziende e dei giochi;
  \item lo storico dei premi dell'industria italiana, in costruzione.
\end{itemize}
Annuncia anche uno \textbf{stream festival} dei giochi italiani che partirà il
lunedì successivo, per quattro giorni, e ricorda il \textbf{marchio collettivo} che
accompagna gli studi italiani nelle fiere internazionali, sostenuto da istituzioni
pubbliche. Ringrazia chi ha costruito il sito con lui e invita a mandare feedback:
il sito serve a \textbf{raccontarsi} quando si gira il mondo\footnote{La
trascrizione si chiude con una frase priva di relazione con l'intervento,
proveniente con ogni probabilità da un altro audio.}.

% ---------------------------------------------------------------------
\section*{Punti chiave del capitolo}
\addcontentsline{toc}{section}{Punti chiave del capitolo}

\begin{itemize}
  \item Il contesto globale: troppi giochi per troppo poca attenzione; anche
        all'estero molti studi arrancano.
  \item Le specificità italiane: industria \textbf{poco densa} e partita tardi,
        confusione tra \textbf{fare un gioco} e \textbf{fare uno studio},
        \textbf{accesso al capitale} difficile, sostegno pubblico tardivo
        (tax credit ``che non basta''; contributi selettivi e automatici della
        legge cinema ancora da attuare).
  \item Dati: 500--600 giochi italiani su Steam, il 42\% con meno di 50
        recensioni ma con gradimento medio sopra l'80\%: la qualità c'è, manca la
        \textbf{scoperta} (e i soldi alla fine); circa 200 aziende e 3.800 persone
        censite.
  \item ``Fare gruppo, non gruppi'': conoscere chi fa che cosa, chiedere,
        sostenere le campagne comuni; scetticismo sì, sospetto no. Come numeri
        \textbf{non siamo ancora un'industria}.
  \item Interventi: fallire presto, pressione anche su regioni e comuni (le
        agevolazioni rendono attrattivi verso i publisher), piccoli fondi per i
        prototipi, giochi brevi: l'obiettivo è \textbf{arrivare al secondo
        gioco}.
  \item Mancano visione editoriale e competenze di mercato e d'impresa, non
        quelle tecniche. Il basso costo del lavoro è un'arma a doppio taglio:
        attrae fondi ma fa \textbf{fuggire il talento}.
  \item Il pregiudizio che i videogiochi non siano cultura sta cedendo; le
        istituzioni ascoltano di più, ma con tempi lenti.
\end{itemize}
